\chapter{Operators, Observables, and the Pauli Matrices}

Good notation is not decoration. In a subject built from abstract vectors and complex numbers, a notation that lets us move symbols without repeatedly translating them into words is part of the machinery. Dirac's bra--ket notation is valuable for precisely this reason: it keeps the geometry of the state space visible while allowing us to calculate with components and matrices.

Our immediate task is mathematical. We shall introduce linear operators, learn how they act on bras as well as kets, and isolate the special class called Hermitian operators. Only then shall we attach the physical words \emph{observable}, \emph{measurement result}, and \emph{probability}. The reward is that the three spin observables emerge as the Pauli matrices rather than being handed to us as objects to memorize.

\section{Orthonormal Bases and Components}

The states form a complex linear vector space. Its ket vectors can be added and multiplied by complex numbers. Associated with every ket \(\ket{A}\) is a bra \(\bra{A}\), obtained by the conjugate-dual operation. A basis is not unique, but once one has been chosen every vector can be described by its components in that basis.

Let
\[
\bigl\{\ket{1},\ket{2},\ldots,\ket{N}\bigr\}
\]
be an orthonormal basis. Orthogonality and normalization are summarized by
\begin{equation}
\braket{j}{i}=\delta_{ji}.
\label{eq:orthonormal-basis}
\end{equation}
Here \(\delta_{ji}\) is one when \(j=i\) and zero otherwise. The dimension \(N\) is the maximum number of mutually orthogonal directions in the space; a complete orthonormal basis contains exactly that many vectors.

Expand an arbitrary ket as
\begin{equation}
\ket{A}=\sum_i \alpha_i\ket{i}.
\label{eq:ket-expansion}
\end{equation}
The coefficients are not mysterious. Multiply from the left by \(\bra{j}\):
\begin{align}
\braket{j}{A}
 &=\sum_i\alpha_i\braket{j}{i} \\
 &=\sum_i\alpha_i\delta_{ji}
 =\alpha_j.
\end{align}
Thus
\begin{equation}
\alpha_i=\braket{i}{A}.
\label{eq:components-inner-products}
\end{equation}
The \(i\)-th component is the projection of \(\ket{A}\) onto the \(i\)-th basis direction.

Substitution of Equation~\eqref{eq:components-inner-products} back into Equation~\eqref{eq:ket-expansion} gives the especially useful form
\begin{equation}
\ket{A}=\sum_i\ket{i}\braket{i}{A}.
\label{eq:ket-completeness-action}
\end{equation}
The sum has done nothing to the vector; it has only resolved the vector into its basis directions and then rebuilt it. We shall use this operation repeatedly. It is usually described as \emph{inserting a complete set of states}, because
\begin{equation}
\sum_i\ket{i}\bra{i}=\mathbf 1.
\label{eq:resolution-identity}
\end{equation}

The corresponding bra expansion is
\begin{equation}
\bra{A}=\sum_i\braket{A}{i}\bra{i}
       =\sum_i\alpha_i^*\bra{i}.
\label{eq:bra-expansion}
\end{equation}
The complex conjugate on \(\alpha_i\) is essential. Reversing a ket into a bra conjugates every numerical coefficient.

\begin{figure}[htbp]
\centering
\includegraphics[width=0.82\textwidth]{lecture_03_figure_04.png}
\caption{The two defining rules of a linear operator, written after the basis review.}
\lectureframe{00:12:44}
\end{figure}

\section{Linear Operators as Machines}

A linear operator \(M\) is a rule that takes an input ket to a unique output ket. A useful picture, due to John Wheeler, is a machine with an input port and an output port: insert \(\ket{A}\), turn the wheel, and receive \(M\ket{A}\). The output is fixed by the input. The reverse need not be unique; two inputs may produce the same output, and an inverse machine need not exist.

The adjective \emph{linear} consists of two rules. For a complex number \(z\),
\begin{align}
M\bigl(z\ket{A}\bigr)&=z\,M\ket{A}, \label{eq:operator-homogeneous}\\
M\bigl(\ket{A}+\ket{B}\bigr)&=M\ket{A}+M\ket{B}.
\label{eq:operator-additive}
\end{align}
A scalar passes through the machine unchanged, and a sum may be processed one term at a time. These rules also imply linearity for any finite complex linear combination.

The same idea applies to ordinary three-dimensional vectors. Quantum mechanics does not invent linear maps. It makes them indispensable because the state space is abstract and complex, and because measurable quantities will be represented by a distinguished family of such maps.

\section{From an Operator to Its Matrix}

Suppose the machine sends \(\ket{A}\) to \(\ket{B}\):
\begin{equation}
M\ket{A}=\ket{B}.
\label{eq:operator-action}
\end{equation}
Choose the same orthonormal basis for input and output and write
\[
\ket{A}=\sum_j\alpha_j\ket{j},
\qquad
\ket{B}=\sum_i\beta_i\ket{i}.
\]
To find the \(i\)-th output component, project Equation~\eqref{eq:operator-action} with \(\bra{i}\):
\begin{equation}
\bra{i}M\ket{A}=\braket{i}{B}=\beta_i.
\end{equation}
Now insert the expansion of \(\ket{A}\):
\begin{align}
\bra{i}M\ket{A}
 &=\sum_j\bra{i}M\ket{j}\braket{j}{A}\\
 &=\sum_j\bra{i}M\ket{j}\alpha_j
 =\beta_i.
\end{align}
The numbers
\begin{equation}
m_{ij}\equiv\bra{i}M\ket{j}
\label{eq:matrix-elements}
\end{equation}
are the \emph{matrix elements} of \(M\) in the chosen basis. Consequently,
\begin{equation}
\boxed{\sum_jm_{ij}\alpha_j=\beta_i.}
\label{eq:component-action}
\end{equation}

\begin{figure}[htbp]
\centering
\includegraphics[width=0.84\textwidth]{lecture_03_figure_02.png}
\caption{The abstract equation \(M\ket{A}=\ket{B}\), projected and resolved into the component rule \(\sum_jm_{ij}\alpha_j=\beta_i\).}
\lectureframe{00:18:30}
\end{figure}

If the space has dimension \(N\), the \(N^2\) numbers \(m_{ij}\) form an \(N\times N\) square matrix. Equation~\eqref{eq:component-action} is precisely
\begin{equation}
\begin{pmatrix}
m_{11}&m_{12}&\cdots&m_{1N}\\
m_{21}&m_{22}&\cdots&m_{2N}\\
\vdots&\vdots&\ddots&\vdots\\
m_{N1}&m_{N2}&\cdots&m_{NN}
\end{pmatrix}
\begin{pmatrix}
\alpha_1\\
\alpha_2\\
\vdots\\
\alpha_N
\end{pmatrix}
=
\begin{pmatrix}
\beta_1\\
\beta_2\\
\vdots\\
\beta_N
\end{pmatrix}.
\label{eq:matrix-vector-action}
\end{equation}
To obtain \(\beta_1\), multiply the first row by the input column and add; repeat with the second row for \(\beta_2\), and so on. The matrix is the concrete version of Wheeler's machine.

\begin{figure}[htbp]
\centering
\includegraphics[width=0.84\textwidth]{lecture_03_figure_03.png}
\caption{The same operator action displayed as a square matrix multiplying the input coefficient column.}
\lectureframe{00:20:01}
\end{figure}

We have gained an algorithm and lost some visible geometry. The abstract statement \(M\ket{A}=\ket{B}\) does not depend on coordinates. The particular entries \(m_{ij}\), \(\alpha_j\), and \(\beta_i\) do depend on the chosen basis. A different orthonormal basis changes all three arrays while leaving the underlying map and its physical predictions unchanged.

\begin{classroomqa}
\audiencequestion{Can the choice of basis determine whether a problem can be solved?}
\lecturerresponse{It can determine whether the solution is easy to see. A well-chosen basis may make an operator diagonal and the calculation almost immediate; a poor basis may hide the structure under unpleasant algebra. It cannot change the result of an experiment. If two bases lead to different physical answers, the calculation is wrong.}
\end{classroomqa}

\section{Bras, Kets, and Anti-Linearity}

\begin{classroomqa}
\audiencequestion{May a linear operator turn a ket into a bra, or a bra into a ket?}
\lecturerresponse{Not an ordinary linear operator of the kind just defined. It maps kets to kets, with complex scalars passing through unchanged. The operation that exchanges kets and bras also complex conjugates their coefficients, so it belongs to a different class.}
\end{classroomqa}

Multiplication by a fixed complex number is a simple linear operator. Complex conjugation is not. If \(K\) denotes conjugation, then
\[
K\bigl(z\ket{A}\bigr)=z^*K\ket{A},
\]
not \(zK\ket{A}\). Such a map is called \emph{anti-linear}. The name does not mean the opposite of linear; it records the conjugation of scalars. Its deeper relation to antiparticles lies beyond the present discussion.

Linear operators can nevertheless be made to act on bras. By convention the operator is written on the right:
\begin{equation}
\bra{A}M.
\end{equation}
The definition is chosen so that a matrix element has no bracket ambiguity:
\begin{equation}
\bra{A}\bigl(M\ket{B}\bigr)
=
\bigl(\bra{A}M\bigr)\ket{B}
\equiv
\bra{A}M\ket{B}.
\label{eq:bracket-independence}
\end{equation}

To see why this is possible, expand both vectors:
\begin{align}
\bra{A}M\ket{B}
 &=\left(\sum_i\alpha_i^*\bra{i}\right)
   M
   \left(\sum_j\beta_j\ket{j}\right)\\
 &=\sum_{ij}\alpha_i^*m_{ij}\beta_j.
\label{eq:matrix-element-double-sum}
\end{align}
This is an ordinary double sum. We may sum over \(j\) first and regard \(M\) as acting on the ket, or sum over \(i\) first and regard it as acting on the bra. The answer is the same. The notation was arranged this way precisely so that unnecessary brackets can be removed.

\section{The Dagger Operation}

There is a second and different bra-side question. Suppose
\[
M\ket{A}=\ket{B}.
\]
What operator turns the corresponding bra \(\bra{A}\) into \(\bra{B}\)? It is not generally \(M\) itself. If \(M\) were multiplication by a complex number \(z\), the ket equation would be \(\ket{B}=z\ket{A}\), while the bra equation would be \(\bra{B}=z^*\bra{A}\).

The required operator is denoted \(M^\dagger\), the \emph{Hermitian conjugate} of \(M\), and is defined by
\begin{equation}
M\ket{A}=\ket{B}
\quad\Longleftrightarrow\quad
\bra{A}M^\dagger=\bra{B}.
\label{eq:dagger-definition}
\end{equation}
The dagger is the operator analogue of complex conjugation.

Its matrix rule follows directly. Let
\[
M\ket{i}=\ket{i'}.
\]
Then
\[
m_{ji}=\bra{j}M\ket{i}=\braket{j}{i'}.
\]
The daggered bra relation is \(\bra{i}M^\dagger=\bra{i'}\), so
\[
(m^\dagger)_{ij}
=\bra{i}M^\dagger\ket{j}
=\braket{i'}{j}.
\]
Reversing an inner product complex conjugates it:
\[
\braket{i'}{j}=\braket{j}{i'}^*.
\]
Therefore
\begin{equation}
\boxed{(m^\dagger)_{ij}=m_{ji}^*.}
\label{eq:dagger-matrix-elements}
\end{equation}

\begin{figure}[htbp]
\centering
\includegraphics[width=0.84\textwidth]{lecture_03_figure_05.png}
\caption{The completed board relation for Hermitian conjugation: interchange the indices and complex conjugate.}
\lectureframe{00:42:23}
\end{figure}

In matrix language we transpose and then complex conjugate:
\begin{equation}
\begin{pmatrix}
a&b\\
c&d
\end{pmatrix}^{\!\dagger}
=
\begin{pmatrix}
a^*&c^*\\
b^*&d^*
\end{pmatrix}.
\label{eq:dagger-two-by-two}
\end{equation}
Transposition reflects the entries across the main diagonal; conjugation replaces every \(i\) by \(-i\). For example,
\begin{equation}
\begin{pmatrix}
2&6+7i\\
4-i&9+i
\end{pmatrix}^{\!\dagger}
=
\begin{pmatrix}
2&4+i\\
6-7i&9-i
\end{pmatrix}.
\label{eq:dagger-numerical-example}
\end{equation}
The diagonal entries remain in their positions but are conjugated, while the two off-diagonal entries exchange positions and are conjugated.

\section{Hermitian Operators}

A complex number is real precisely when it equals its own complex conjugate. The operator analogue is immediate.

\begin{definition}
An operator \(M\) is \emph{Hermitian} when
\begin{equation}
M=M^\dagger.
\label{eq:hermitian-definition}
\end{equation}
Equivalently, in every orthonormal basis,
\begin{equation}
m_{ij}=m_{ji}^*.
\label{eq:hermitian-components}
\end{equation}
\end{definition}

Set \(i=j\) in Equation~\eqref{eq:hermitian-components}. Then
\[
m_{ii}=m_{ii}^*,
\]
so every diagonal entry is real. Away from the diagonal the entries occur in conjugate-reflected pairs. If one entry is \(4-i\), its reflected partner is \(4+i\); if one is the real number \(5\), so is its partner.

\begin{figure}[htbp]
\centering
\includegraphics[width=0.78\textwidth]{lecture_03_figure_06.png}
\caption{A completed Hermitian-matrix example: real entries on the diagonal and complex-conjugate pairs reflected across it.}
\lectureframe{00:45:43}
\end{figure}

Hermitian operators matter because they will represent measurable quantities. Before making that identification, we need to understand their eigendirections.

\section{Eigenvectors and Eigendirections}

A first encounter with eigenvectors can make this part feel faster than the physics around it. The ideas are standard linear algebra rather than a new physical hypothesis; if they are unfamiliar, it is worth reviewing the definitions quietly and then returning to the argument.

A general operator changes the direction of a vector. A rotation in ordinary three-dimensional space, for example, turns almost every vector away from itself. The rotation axis is exceptional: a vector along that axis returns in the same direction. This is the model for an eigenvector.

A nontrivial rotation of a real two-dimensional plane illustrates the opposite possibility: there is no real direction left fixed up to a real scale. Thus a generic operator on a restricted real space need not have a real eigenvector. The Hermitian operators on the complex state spaces used in quantum mechanics have much stronger properties.

An eigenvector of \(M\) satisfies
\begin{equation}
M\ket{i}=\lambda_i\ket{i}.
\label{eq:eigenvalue-equation}
\end{equation}
The number \(\lambda_i\) is its eigenvalue. The eigenvector belongs to the operator; one does not speak of an eigenvector without specifying the map whose direction it preserves.

Its length and phase are not fixed by Equation~\eqref{eq:eigenvalue-equation}. For every nonzero complex number \(z\),
\[
M\bigl(z\ket{i}\bigr)
=zM\ket{i}
=\lambda_i\bigl(z\ket{i}\bigr).
\]
It is therefore more accurate to speak of an \emph{eigendirection}: every nonzero multiple of \(\ket{i}\) describes the same direction and carries the same eigenvalue. We usually choose a normalized representative.

The central spectral fact is this: a Hermitian operator on an \(N\)-dimensional complex vector space has a complete orthonormal basis of \(N\) eigenvectors. Two parts of that statement can be seen with very little algebra.

\subsection{Hermitian Eigenvalues Are Real}

Assume \(M=M^\dagger\) and
\[
M\ket{i}=\lambda_i\ket{i}.
\]
Take the Hermitian conjugate of the whole equation:
\begin{equation}
\bra{i}M^\dagger=\lambda_i^*\bra{i}.
\label{eq:eigen-bra-equation}
\end{equation}
Since \(M^\dagger=M\), multiply the ket equation by \(\bra{i}\) and the bra equation by \(\ket{i}\):
\begin{align}
\bra{i}M\ket{i}&=\lambda_i\braket{i}{i},\\
\bra{i}M\ket{i}&=\lambda_i^*\braket{i}{i}.
\end{align}
The same nonzero norm multiplies both numbers, so
\begin{equation}
\lambda_i=\lambda_i^*.
\end{equation}
Every eigenvalue of a Hermitian operator is real. This is the first hint that such an operator can represent the real number displayed by a measuring apparatus.

\subsection{Different Eigenvalues Give Orthogonal Eigenvectors}

Let
\[
M\ket{i}=\lambda_i\ket{i},
\qquad
M\ket{j}=\lambda_j\ket{j},
\]
where \(M\) is Hermitian. Dagger the second equation. Its eigenvalue is already known to be real:
\[
\bra{j}M=\lambda_j\bra{j}.
\]
Now multiply the first equation by \(\bra{j}\) and the bra equation by \(\ket{i}\):
\begin{align}
\bra{j}M\ket{i}&=\lambda_i\braket{j}{i},\\
\bra{j}M\ket{i}&=\lambda_j\braket{j}{i}.
\end{align}
Subtracting gives
\begin{equation}
\boxed{(\lambda_i-\lambda_j)\braket{j}{i}=0.}
\label{eq:hermitian-orthogonality}
\end{equation}
If \(\lambda_i\ne\lambda_j\), then
\[
\braket{j}{i}=0.
\]
The eigendirections are orthogonal.

\begin{figure}[htbp]
\centering
\includegraphics[width=0.82\textwidth]{lecture_03_figure_07.png}
\caption{The decisive line in the orthogonality proof: unequal eigenvalues force the inner product to vanish.}
\lectureframe{00:59:00}
\end{figure}

When two eigenvalues happen to be equal, Equation~\eqref{eq:hermitian-orthogonality} imposes no condition within that degenerate eigenspace. One may choose an orthonormal basis there. Together with the remaining eigendirections, those vectors form the complete orthonormal eigenbasis promised by the spectral theorem.

Orthogonality also has a physical reading. Two orthogonal states can be distinguished without ambiguity by a suitable measurement. Eigenstates with different eigenvalues must be distinguishable because the apparatus reports different real numbers for them. The algebra and the operational meaning are already pointing to the same correspondence.

\section{The Measurement Postulates}

We now have enough formalism to state the basic quantum-mechanical correspondences. They may sound foreign when first encountered. Their meaning becomes concrete only after they are used, so let us state them cleanly and then return at once to the spin experiment.

A single measurement produces a real number. One may package two compatible real readings into a complex number, but that package still contains two measurements. The elementary output of one detector is real. Hermitian operators are therefore the natural candidates: their eigenvalues are real, and their eigenvectors form a complete orthonormal basis.

\begin{enumerate}
\item A measurable quantity, or \emph{observable}, is represented by a Hermitian operator \(M\).
\item The possible results of measuring \(M\) are its eigenvalues \(\lambda_i\).
\item A state in which \(M\) has the definite value \(\lambda_i\) is an eigenstate \(\ket{i}\) satisfying
\[
M\ket{i}=\lambda_i\ket{i}.
\]
\item If the system is in a general normalized state \(\ket{A}\), then the probability of obtaining \(\lambda_i\) is
\begin{equation}
P(\lambda_i\mid A)
=\left|\braket{i}{A}\right|^2
=\braket{i}{A}\braket{A}{i}.
\label{eq:born-rule}
\end{equation}
\end{enumerate}

This list is a convenient starting point, not necessarily a minimal independent axiom system; with a more economical formulation, some statements can be derived from others. Keeping the correspondences separate makes their physical roles easier to see.

The complex number \(\braket{i}{A}\) is a \emph{probability amplitude}, not a probability. Multiplication by its complex conjugate produces a real, nonnegative number. For a spin component, the spectrum has only two values, \(+1\) and \(-1\); a more general observable may have many eigenvalues.

\begin{figure}[htbp]
\centering
\includegraphics[width=0.82\textwidth]{lecture_03_figure_08.png}
\caption{The board connects definite-result states with eigenvectors and the probability of \(\lambda_i\) with the squared projection onto \(\ket{i}\).}
\lectureframe{01:11:29}
\end{figure}

\begin{classroomqa}
\audiencequestion{What exactly is written next to \(P\) in the probability rule?}
\lecturerresponse{It is \(\lambda_i\): \(P(\lambda_i\mid A)\) means the probability that a measurement of \(M\), when the system is prepared in \(\ket{A}\), returns the particular eigenvalue \(\lambda_i\). The right-hand side is the squared magnitude of the projection onto that eigenvector.}
\end{classroomqa}

\subsection{Normalization and Completeness}

The probabilities of all mutually exclusive outcomes must add to one:
\begin{equation}
\sum_i\left|\braket{i}{A}\right|^2=1.
\label{eq:probability-normalization}
\end{equation}
Insert the complete eigenbasis:
\begin{align}
\sum_i\left|\braket{i}{A}\right|^2
 &=\sum_i\braket{A}{i}\braket{i}{A}\\
 &=\bra{A}\left(\sum_i\ket{i}\bra{i}\right)\ket{A}\\
 &=\braket{A}{A}.
\end{align}
Hence an allowable state vector is normalized,
\begin{equation}
\braket{A}{A}=1.
\label{eq:state-normalization}
\end{equation}
If \(\ket{A}=\sum_i\alpha_i\ket{i}\), the same condition is
\begin{equation}
\sum_i\alpha_i^*\alpha_i=1.
\label{eq:component-normalization}
\end{equation}
This is one real constraint on the complex components.

\begin{classroomqa}
\audiencequestion{Must every probability be positive?}
\lecturerresponse{Yes. Each probability is a number times its own complex conjugate, so it is nonnegative. Normalization then makes the sum of all such numbers equal to one.}
\end{classroomqa}

There is another redundancy. Multiplying every component by the same phase,
\begin{equation}
\ket{A}\longmapsto e^{i\theta}\ket{A},
\label{eq:global-phase}
\end{equation}
multiplies every amplitude by \(e^{i\theta}\) and leaves every squared magnitude unchanged. No measurement can detect that common phase. Relative phases between components remain physical; only a phase multiplying the whole state is redundant.

Thus two real parameters in an arbitrary complex vector do not specify the physical state: its magnitude is fixed by normalization, and its overall phase is unobservable. This is the same parameter count encountered for a single spin, now derived from the general probability rule.

\subsection{Measurement Also Prepares}

One further postulate is easier to say than to compress into a short blackboard line. If a measurement returns \(\lambda_i\), then immediately afterward the system is in the corresponding eigenstate \(\ket{i}\). The apparatus both records a result and prepares a state in which a repetition of the same measurement has that definite result.

\begin{classroomqa}
\audiencequestion{Could that measurement-preparation statement be said once more?}
\lecturerresponse{Before the measurement, several eigenvalues may occur with the probabilities in Equation~\eqref{eq:born-rule}. Once the apparatus actually reports \(\lambda_i\), the system is left in the eigenstate belonging to \(\lambda_i\). A second immediate measurement of the same observable therefore returns the same value.}
\end{classroomqa}

\begin{classroomqa}
\audiencequestion{Do these principles already tell us how the state evolves with time?}
\lecturerresponse{Not yet. We have specified what a state is, how observables are represented, and what a measurement does. The law governing change with time is additional structure and will be introduced later. This is analogous to first defining classical phase space before writing equations of motion.}
\end{classroomqa}

\begin{classroomqa}
\audiencequestion{Earlier the average spin along a tilted axis involved \(\cos\theta\). Should the new squared-amplitude rule make that \(\cos^2\theta\)?}
\lecturerresponse{The probability and the average are different quantities. The average combines the probabilities of the \(+1\) and \(-1\) outcomes with their signs. Once the arbitrary-direction spinor is constructed, the squared amplitudes reproduce the earlier cosine average. That calculation is the next application, not something to guess by changing the old formula.}
\end{classroomqa}

\section{Spin as the First Application}

The three measurable spin components are represented by Hermitian operators
\[
\sigma_x,\qquad \sigma_y,\qquad \sigma_z.
\]
The analyzer experiments already supplied their possible values and their definite-value states:
\[
\begin{array}{c|c|c}
\text{observable}&\text{eigenvalue }+1&\text{eigenvalue }-1\\ \hline
\sigma_z&\ket{u}&\ket{d}\\
\sigma_x&\ket{r}&\ket{l}\\
\sigma_y&\ket{\mathrm{in}}&\ket{\mathrm{out}}
\end{array}
\]
This information is enough to reconstruct the operators as matrices.

\subsection{The z Component in Its Eigenbasis}

Choose the up--down basis. Definite \(z\)-spin means
\begin{equation}
\sigma_z\ket{u}=+\ket{u},
\qquad
\sigma_z\ket{d}=-\ket{d}.
\label{eq:sigma-z-eigenvectors}
\end{equation}
The two eigenvalues are different, so Hermiticity implies \(\braket{u}{d}=0\).

Compute the four matrix elements:
\begin{align}
\bra{u}\sigma_z\ket{u}&=+1,&
\bra{u}\sigma_z\ket{d}&=-\braket{u}{d}=0,\\
\bra{d}\sigma_z\ket{u}&=\braket{d}{u}=0,&
\bra{d}\sigma_z\ket{d}&=-1.
\end{align}
Therefore
\begin{equation}
\boxed{
\sigma_z=
\begin{pmatrix}
1&0\\
0&-1
\end{pmatrix}_{\{u,d\}}.}
\label{eq:sigma-z}
\end{equation}
The coordinate columns are
\[
\ket{u}\widehat{=}\begin{pmatrix}1\\0\end{pmatrix},
\qquad
\ket{d}\widehat{=}\begin{pmatrix}0\\1\end{pmatrix}.
\]
Multiplication by Equation~\eqref{eq:sigma-z} returns the first column unchanged and the second with a minus sign, exactly as the eigenvalue equations require.

\begin{figure}[htbp]
\centering
\includegraphics[width=0.82\textwidth]{lecture_03_figure_09.png}
\caption{The \(\sigma_z\) eigenvalue equations and the resulting diagonal matrix in the up--down basis.}
\lectureframe{01:26:56}
\end{figure}

\subsection{The x Component Without Changing Basis}

In the right--left basis, \(\sigma_x\) would have the same diagonal form. That would be easy but uninformative. Keep the up--down basis and use
\begin{equation}
\sigma_x\ket{r}=+\ket{r},
\qquad
\sigma_x\ket{l}=-\ket{l}.
\label{eq:sigma-x-eigenvectors}
\end{equation}
The two bases are related by
\begin{align}
\ket{r}&=\frac{\ket{u}+\ket{d}}{\sqrt2},&
\ket{l}&=\frac{\ket{u}-\ket{d}}{\sqrt2},\\
\ket{u}&=\frac{\ket{r}+\ket{l}}{\sqrt2},&
\ket{d}&=\frac{\ket{r}-\ket{l}}{\sqrt2}.
\label{eq:x-z-basis-relations}
\end{align}

Begin with the upper diagonal entry:
\begin{align}
\bra{u}\sigma_x\ket{u}
 &=\frac12(\bra{r}+\bra{l})\sigma_x(\ket{r}+\ket{l})\\
 &=\frac12(\bra{r}+\bra{l})(\ket{r}-\ket{l})\\
 &=0.
\end{align}
The right and left states are orthonormal, so the two diagonal terms cancel and the cross terms vanish. The same calculation with \(\ket{d}\) gives
\[
\bra{d}\sigma_x\ket{d}=0.
\]
At this point two zeros do not mean the operator vanishes; they tell us to look off the diagonal.

For the upper-right entry,
\begin{align}
\bra{u}\sigma_x\ket{d}
 &=\frac12(\bra{r}+\bra{l})\sigma_x(\ket{r}-\ket{l})\\
 &=\frac12(\bra{r}+\bra{l})(\ket{r}+\ket{l})\\
 &=1.
\end{align}
Hermiticity, or a direct repetition of the calculation, gives
\[
\bra{d}\sigma_x\ket{u}=1.
\]
Thus
\begin{equation}
\boxed{
\sigma_x=
\begin{pmatrix}
0&1\\
1&0
\end{pmatrix}_{\{u,d\}}.}
\label{eq:sigma-x}
\end{equation}
The same abstract operator is diagonal in its own eigenbasis and off-diagonal in the \(z\)-basis. This is a concrete instance of basis dependence without physical ambiguity.

\subsection{The y Component and the Relative Phase}

Choose the normalized in--out convention
\begin{equation}
\ket{\mathrm{in}}=\frac{\ket{u}+i\ket{d}}{\sqrt2},
\qquad
\ket{\mathrm{out}}=\frac{\ket{u}-i\ket{d}}{\sqrt2},
\label{eq:in-out-states}
\end{equation}
with
\[
\sigma_y\ket{\mathrm{in}}=+\ket{\mathrm{in}},
\qquad
\sigma_y\ket{\mathrm{out}}=-\ket{\mathrm{out}}.
\]
Interchanging the labels or reversing the \(y\)-axis changes a sign convention, not the underlying construction. With Equation~\eqref{eq:in-out-states}, the unique Hermitian matrix with these eigenvectors and eigenvalues is
\begin{equation}
\boxed{
\sigma_y=
\begin{pmatrix}
0&-i\\
i&0
\end{pmatrix}_{\{u,d\}}.}
\label{eq:sigma-y}
\end{equation}
It is easy to verify rather than rederive:
\begin{align}
\begin{pmatrix}0&-i\\i&0\end{pmatrix}
\frac1{\sqrt2}\begin{pmatrix}1\\i\end{pmatrix}
&=\frac1{\sqrt2}\begin{pmatrix}1\\i\end{pmatrix},\\
\begin{pmatrix}0&-i\\i&0\end{pmatrix}
\frac1{\sqrt2}\begin{pmatrix}1\\-i\end{pmatrix}
&=-\frac1{\sqrt2}\begin{pmatrix}1\\-i\end{pmatrix}.
\end{align}
The factors of \(1/\sqrt2\) normalize the eigenvectors but do not affect their eigendirections.

\begin{figure}[htbp]
\centering
\includegraphics[width=0.82\textwidth]{lecture_03_figure_10.png}
\caption{The completed board set of Pauli matrices in the up--down basis.}
\lectureframe{01:38:11}
\end{figure}

\begin{classroomqa}
\audiencequestion{Should one sign in the in--out vectors be changed to make them orthogonal?}
\lecturerresponse{No. The missing step is complex conjugation when the ket is turned into a bra:
\[
\bra{\mathrm{in}}=\frac{\bra{u}-i\bra{d}}{\sqrt2}.
\]
Therefore
\[
\braket{\mathrm{in}}{\mathrm{out}}
=\frac12(\bra{u}-i\bra{d})(\ket{u}-i\ket{d})
=\frac12\bigl(1+(-i)(-i)\bigr)=0.
\]
Forgetting that conjugation is precisely the common mistake that bra--ket notation is meant to expose.}
\end{classroomqa}

We have obtained
\begin{equation}
\sigma_x=\begin{pmatrix}0&1\\1&0\end{pmatrix},
\qquad
\sigma_y=\begin{pmatrix}0&-i\\i&0\end{pmatrix},
\qquad
\sigma_z=\begin{pmatrix}1&0\\0&-1\end{pmatrix}.
\label{eq:pauli-triple}
\end{equation}
They are Hermitian, have eigenvalues \(\pm1\), and possess exactly the definite-spin eigenvectors demanded by the analyzer experiments.

The next problem is now well posed. Given a spin prepared along an arbitrary direction, construct the corresponding state and use squared amplitudes to calculate the probabilities for \(\sigma_z=+1\) and \(\sigma_z=-1\). The formalism already contains enough information; carrying out that calculation is the natural next step.
